\chapter{1994 Nobel Prize in Economics}
\textbf{Laureates: }John Nash (USA), Reinhard Selten (Germany), John Harsanyi (Hungary/USA)

\section*{Reason for Award}
For their pioneering analysis of equilibria in the theory of non-cooperative games

\section{What They Did}
John Nash introduced the Nash equilibrium concept, defining a situation in which
each player's strategy is a best response to the strategies of all other players.
In a Nash equilibrium, no player can improve their payoff by unilaterally
deviating from their strategy. This became a foundational solution concept for
non-cooperative games and a rigorous framework for analyzing strategic
interaction. Nash proved that every finite game has at least one equilibrium in
mixed or pure strategies, using a fixed-point argument.

Reinhard Selten refined the equilibrium concept by introducing subgame
perfection. He identified that some Nash equilibria rely on non-credible threats
in dynamic games. Subgame perfection eliminates some of these equilibria by
requiring that strategies constitute Nash equilibria in every subgame of the
original game. This refinement requires threats and promises to be sequentially
rational in proper subgames, not merely on the equilibrium path.

John Harsanyi extended game theory to situations of incomplete information,
where players possess private knowledge about their own types, preferences, or
payoffs. Harsanyi introduced the concept of types, where each player's type
represents private information. Using type spaces and, in many models, a common
prior, he showed how players could form beliefs about others' types and analyze
strategic choices under asymmetric information. Bayesian Nash equilibrium
became a standard tool for analyzing auctions, bargaining, and other games with
private information.

Together, these three laureates supplied central solution concepts and methods
that greatly expanded the usefulness of non-cooperative game theory. Game-
theoretic methods are now used in economics, political science, biology,
computer science, philosophy, and international relations.

\subsection*{Nash Equilibrium}
The Nash equilibrium represents a stable state of play where no player has an
incentive to deviate unilaterally. Nash's proof relied on fixed-point theorems,
showing that best-response correspondences have equilibria under conditions
such as finite strategy sets or appropriate compactness and continuity.
The concept applies to both simultaneous and sequential games, though it does
not always select the most plausible outcome in dynamic settings.

\subsection*{Subgame Perfection}
Selten's subgame perfection requires that strategies form Nash equilibria in
every proper subgame. This refinement eliminates equilibria sustained by
non-credible threats. For example, in a sequential entry game, a monopolist's
threat to fight entry may not be credible if fighting is not optimal once entry
has occurred. Subgame perfection requires sequential rationality in every
proper subgame; in games with appropriate perfect-information structure, this
can be checked by backward induction. It does not guarantee a unique or
empirically correct prediction.

\subsection*{Bayesian Games}
Harsanyi's innovation was to represent games of incomplete information as
games of imperfect information by introducing nature as a device that selects
players' types. This transformation allows beliefs and Nash equilibrium
concepts to be used when players have private information. Harsanyi's approach
became a foundation for analyzing auctions, principal--agent problems, and
signaling games.

\section{Impact on Economic Thinking}
Nash, Selten, and Harsanyi collectively transformed the analysis of strategic
decision-making by providing rigorous solution concepts. Equilibrium concepts
help researchers analyze situations where each agent's optimal choice depends
on the choices and information of others, although equilibria may be multiple
or may not predict observed behavior without additional assumptions.

Self-enforcing strategies allow researchers to analyze complex strategic
environments systematically. Game theory became an important tool for studying
oligopoly behavior, auctions, mechanism design, coalition formation,
bargaining, negotiation, and conflict. It also found applications in
evolutionary biology, where stable strategies can help explain population
dynamics.

Mechanism design, principal--agent models, signaling, and screening build on
game-theoretic analysis of incentives and private information. Nash-equilibrium
refinements address some multiplicity problems, with subgame perfection
eliminating equilibria supported by certain non-credible threats and Harsanyi's
type-space formulation handling incomplete information.

Strategic analysis permeates modern economic research across industrial
organization, macroeconomics, finance, and public economics. Game theory has
become the microfoundations of much of modern economics, providing a unified
language for analyzing interactive decision making. Its mathematical elegance
paired with profound practical applications stands among the greatest achievements
of twentieth-century economic thought.

\subsection*{Industrial Organization}
Game theory transformed industrial organization by providing tools for analyzing
oligopolistic competition, collusion, entry deterrence, and product differentiation.
Cournot, Bertrand, and Stackelberg models can be formulated as games with
corresponding equilibrium concepts, including Nash or subgame-perfect
equilibrium depending on the setting. Modern industrial-organization research
relies heavily on game-theoretic methods.

\subsection*{Auction Theory}
Harsanyi's framework for incomplete information became important in auction
theory, where researchers analyze bidding strategies, information, and revenue.
Vickrey, sealed-bid, and ascending auctions can be studied as Bayesian games.
Auction design has practical applications in spectrum licensing, emissions
trading, and procurement.

\section{Modern Perspectives Today}
Mechanism design uses game-theoretic foundations to design incentives and rules
for specific objectives. Auction formats for spectrum allocation, as well as
matching and school-assignment mechanisms, draw on the broader game-theory and
mechanism-design literature that followed these contributions.
Algorithmic game theory combines computational complexity with game theoretic
solution concepts.

Evolutionary game theory analyzes population dynamics and stable strategies,
with applications to biology and economics. Behavioral game theory incorporates
psychological realism into equilibrium analysis, relaxing the strict rationality
assumptions. Neuroeconomics investigates the neural correlates of strategic
decision making.

Algorithmic mechanism design addresses computational constraints in implementing
mechanisms. Game theory also informs multi-agent systems, artificial
intelligence, autonomous-vehicle coordination, digital platforms, and security
analysis, although these applications often require models beyond the original
equilibrium frameworks.

International climate agreements are sometimes modeled as repeated or dynamic
games, while common-pool resources are studied using strategic-interaction
models. Extending the contributions of Nash, Selten, and Harsanyi, game theory
continues to develop through computational, behavioral, spatial, and networked
approaches.

\subsection*{Key References}
\begin{itemize}
    \item Nash, J. F. (1950). ``Equilibrium Points in $n$-Person Games.'' \textit{Proceedings of the National Academy of Sciences}, 36(1), 48--49.
    \item Harsanyi, J. C. (1967--1968). ``Games with Incomplete Information Played by `Bayesian' Players,'' Parts I--III. \textit{Management Science}, 14.
    \item Selten, R. (1975). ``Reexamination of the Perfectness Concept for Equilibrium Points in Extensive Games.'' \textit{International Journal of Game Theory}, 4, 25--55.
\end{itemize}
